% GENERATED FILE. Edit canonical lessons, then run npm run generate:books.
% canonical-source-hash: fnv1a64:bd966224d315423d
% canonical-lessons: MR-C258-halli, MR-C258-satak, MR-C258-yug, MR-C258-ratrabhar, MR-C258-divasbhar

\chapter{Nowadays, Century, Era, All night, All day}
\label{ch:mr-nowadays-century-era-all-night-all-day}
\hlchaptermodality{258}

\begin{chapteropening}
\textbf{By the end of this chapter:} \emph{I can say nowadays, a century, an era, all night and all day.}

Complete the last lesson of chapter 258: I can say nowadays, a century, an era, all night and all day.
\end{chapteropening}

\section[\texorpdfstring{\mr{हल्ली} (hallī)}{हल्ली (hallī)}]{\mr{हल्ली} (hallī) — nowadays}

\label{lesson:MR-C258-halli}



\begin{quote}
Before the new one: say the Marathi for trust, then the Marathi for courage.
\end{quote}

\subsection*{You'll want to know: \mr{हल्ली}}
\textbf{\mr{हल्ली}} — \emph{hallī} — \textquotedblleft{}nowadays\textquotedblright{}.

\begin{cousinweb}
The origin of \textbf{\mr{हल्ली}} (\emph{hallī}) is not settled. Neither CDIAL nor DEDR, searched by form and by meaning, has an entry for it.
\end{cousinweb}

\begin{grammarlens}[title={What you've built}]
One more word for this chapter.
\end{grammarlens}

\subsection*{Your turn}
\begin{itemize}
\raggedright
  \item \emph{Say it:} \emph{hallī}
  \item \emph{Say it:} \emph{hallī}, once more
  \item \emph{Say it:} say \emph{hallī}
  \item \emph{Recall:} say the Marathi for trust, then the Marathi for courage, then say \emph{hallī} again
\end{itemize}

\begin{tcolorbox}[breakable,colback=teal!4,colframe=teal!35!black,title={Before you move on}]
What does \mr{हल्ली} mean? (\textquotedblleft{}Nowadays\textquotedblright{}.) Say it once more.
\end{tcolorbox}
\section[\texorpdfstring{\mr{शतक} (śatak)}{शतक (śatak)}]{\mr{शतक} (śatak) — a century}

\label{lesson:MR-C258-satak}



\begin{quote}
Before the new one: say the Marathi for courage, then the Marathi for nowadays.
\end{quote}

\subsection*{You'll want to know: \mr{शतक}}
\textbf{\mr{शतक}} — \emph{śatak} — \textquotedblleft{}a century\textquotedblright{}.

\begin{cousinweb}
\textbf{\mr{शतक}} (\emph{śatak}) is Sanskrit \emph{śataka}, \textquotedblleft{}consisting of a hundred\textquotedblright{}, a tatsama built on \emph{śata}, \textquotedblleft{}hundred\textquotedblright{}. CDIAL has no entry for it.
\end{cousinweb}

\begin{grammarlens}[title={What you've built}]
One more word for this chapter.
\end{grammarlens}

\subsection*{Your turn}
\begin{itemize}
\raggedright
  \item \emph{Say it:} \emph{śatak}
  \item \emph{Say it:} \emph{śatak}, once more
  \item \emph{Say it:} say \emph{śatak}
  \item \emph{Recall:} say the Marathi for courage, then the Marathi for nowadays, then say \emph{śatak} again
\end{itemize}

\begin{tcolorbox}[breakable,colback=teal!4,colframe=teal!35!black,title={Before you move on}]
What does \mr{शतक} mean? (\textquotedblleft{}A century\textquotedblright{}.) Say it once more.
\end{tcolorbox}
\section[\texorpdfstring{\mr{युग} (yug)}{युग (yug)}]{\mr{युग} (yug) — an era, an age}

\label{lesson:MR-C258-yug}



\begin{quote}
Before the new one: say the Marathi for nowadays, then the Marathi for a century.
\end{quote}

\subsection*{You'll want to know: \mr{युग}}
\textbf{\mr{युग}} — \emph{yug} — \textquotedblleft{}an era, an age\textquotedblright{}.

\begin{cousinweb}
\textbf{\mr{युग}} (\emph{yug}) is Sanskrit \emph{yuga}, \textquotedblleft{}yoke; couple\textquotedblright{}, taken over whole as a learned loan, a tatsama (CDIAL 10482). Sanskrit \emph{yuga} also meant \textquotedblleft{}an age of the world\textquotedblright{}. CDIAL lists the inherited Marathi form as \emph{jūṁ}, \textquotedblleft{}yoke\textquotedblright{}.
\end{cousinweb}

\begin{grammarlens}[title={What you've built}]
One more word for this chapter.
\end{grammarlens}

\subsection*{Your turn}
\begin{itemize}
\raggedright
  \item \emph{Say it:} \emph{yug}
  \item \emph{Say it:} \emph{yug}, once more
  \item \emph{Say it:} say \emph{yug}
  \item \emph{Recall:} say the Marathi for nowadays, then the Marathi for a century, then say \emph{yug} again
\end{itemize}

\begin{tcolorbox}[breakable,colback=teal!4,colframe=teal!35!black,title={Before you move on}]
What does \mr{युग} mean? (\textquotedblleft{}An era, an age\textquotedblright{}.) Say it once more.
\end{tcolorbox}
\section[\texorpdfstring{\mr{रात्रभर} (rātrabhar)}{रात्रभर (rātrabhar)}]{\mr{रात्रभर} (rātrabhar) — all night}

\label{lesson:MR-C258-ratrabhar}



\begin{quote}
Before the new one: say the Marathi for a century, then the Marathi for an era.
\end{quote}

\subsection*{You'll want to know: \mr{रात्रभर}}
\textbf{\mr{रात्रभर}} — \emph{rātrabhar} — \textquotedblleft{}all night\textquotedblright{}.

\begin{cousinweb}
\textbf{\mr{रात्रभर}} (\emph{rātrabhar}) is Sanskrit \emph{rātra}, \textquotedblleft{}night\textquotedblright{}, with \emph{-bhar}, \textquotedblleft{}full\textquotedblright{}, from \emph{bharṇẽ}, \textquotedblleft{}to fill\textquotedblright{} (CDIAL 9397): \textquotedblleft{}the night full\textquotedblright{}.
\end{cousinweb}

\begin{grammarlens}[title={What you've built}]
One more word for this chapter.
\end{grammarlens}

\subsection*{Your turn}
\begin{itemize}
\raggedright
  \item \emph{Say it:} \emph{rātrabhar}
  \item \emph{Say it:} \emph{rātrabhar}, once more
  \item \emph{Say it:} say \emph{rātrabhar}
  \item \emph{Recall:} say the Marathi for a century, then the Marathi for an era, then say \emph{rātrabhar} again
\end{itemize}

\begin{tcolorbox}[breakable,colback=teal!4,colframe=teal!35!black,title={Before you move on}]
What does \mr{रात्रभर} mean? (\textquotedblleft{}All night\textquotedblright{}.) Say it once more.
\end{tcolorbox}
\section[\texorpdfstring{\mr{दिवसभर} (divasbhar)}{दिवसभर (divasbhar)}]{\mr{दिवसभर} (divasbhar) — all day}

\label{lesson:MR-C258-divasbhar}



\begin{quote}
Before the new one: say the Marathi for an era, then the Marathi for all night.
\end{quote}

\subsection*{You'll want to know: \mr{दिवसभर}}
\textbf{\mr{दिवसभर}} — \emph{divasbhar} — \textquotedblleft{}all day\textquotedblright{}.

\begin{cousinweb}
\textbf{\mr{दिवसभर}} (\emph{divasbhar}) is \emph{divas}, \textquotedblleft{}day\textquotedblright{}, from Sanskrit \emph{divasa} (CDIAL 6333), with \emph{-bhar}, \textquotedblleft{}full\textquotedblright{}, from \emph{bharṇẽ}, \textquotedblleft{}to fill\textquotedblright{} (CDIAL 9397).
\end{cousinweb}

\begin{grammarlens}[title={What you've built}]
One more word for this chapter.
\end{grammarlens}

\subsection*{Your turn}
\begin{itemize}
\raggedright
  \item \emph{Say it:} \emph{divasbhar}
  \item \emph{Say it:} \emph{divasbhar}, once more
  \item \emph{Say it:} say \emph{divasbhar}
  \item \emph{Recall:} say the Marathi for an era, then the Marathi for all night, then say \emph{divasbhar} again
\end{itemize}

\begin{tcolorbox}[breakable,colback=teal!4,colframe=teal!35!black,title={Before you move on}]
What does \mr{दिवसभर} mean? (\textquotedblleft{}All day\textquotedblright{}.) Say it once more.
\end{tcolorbox}
